% ============================================================================
% BOLUM 1: OZET VE GIRIS
% ============================================================================
% Kurumsal Temizlik Takip ve Yonetim Sistemi (KTTYS)
% Muhammet Anil YAGIZ - 213255046
% Kirikkale Universitesi - Bilgisayar Muhendisligi Bolumu
% ============================================================================

\chapter{OZET VE GIRIS}
\label{chap:ozet}

% ============================================================================
% 1.1 OZET
% ============================================================================
\section{Ozet}
\label{sec:ozet}

Gunumuzde kurumsal yonetim sistemleri, organizasyonlarin verimli ve etkin bir 
sekilde isleyisini surdurebilmesi icin kritik bir oneme sahiptir. Ozellikle 
buyuk olcekli kurumlarda temizlik hizmetlerinin planlanmasi, takibi ve 
denetlenmesi ciddi bir koordinasyon sorunu olusturmaktadir. Bu calismada, 
Kirikkale Universitesi kampusundeki temizlik operasyonlarini dijitallestirmek 
ve optimize etmek amaciyla \textbf{Kurumsal Temizlik Takip ve Yonetim Sistemi 
(KTTYS)} gelistirilmistir.

Gelistirilen sistem, modern web teknolojileri kullanilarak tasarlanmis olup, 
Progressive Web Application (PWA) mimarisi uzerine insa edilmistir. Sistem, 
Node.js tabanli Express.js framework'u ile RESTful API mimarisi kullanilarak 
backend tarafinda, SQLite veritabani ile veri kaliciligi saglanarak ve 
responsive HTML5/CSS3/JavaScript ile frontend tarafinda tam entegre bir 
cozum sunmaktadir.

KTTYS, cok katmanli bir kullanici yetkilendirme sistemi sunarak, Super Admin, 
Fakulte Yoneticisi, Denetleyici (Supervisor) ve Temizlik Personeli olmak uzere 
dort farkli kullanici rolunu desteklemektedir. Her bir rol, sistem uzerinde 
belirli yetkilere sahip olup, hiyerarsik bir yonetim yapisi olusturmaktadir. 
Bu yapi sayesinde, universite genelindeki temizlik operasyonlari merkezi 
olarak yonetilebilmekte, ancak fakulte bazinda desantralize bir kontrol 
mekanizmasi da saglanmaktadir.

Sistemin temel ozellikleri arasinda haftalik temizlik programlarinin 
olusturulmasi, gunluk gorevlerin otomatik atanmasi, gorev tamamlama ve onay 
sureclerinin takibi, performans raporlarinin olusturulmasi ve resmi tatil 
gunlerinin yonetimi yer almaktadir. Ayrica sistem, Turkiye'deki resmi tatiller 
ve dini bayramlari otomatik olarak taniyarak, bu gunlerde gorev atamasi 
yapilmamasini saglamaktadir.

Calisma sonucunda, manuel olarak yurutulen temizlik takip sureclerinin 
dijitallestirilmesiyle zaman ve kaynak tasarrufu saglanmistir. Sistem, 
kullanici dostu arayuzu ve mobil uyumlu tasarimi sayesinde tum paydaslar 
tarafindan kolaylikla kullanilabilmektedir. Gerceklestirilen testler, 
sistemin yuksek performans ve guvenilirlik standartlarini karsiladigini 
gostermektedir.

\textbf{Anahtar Kelimeler:} Temizlik Takip Sistemi, Kurumsal Yonetim, 
Progressive Web Application, Node.js, Express.js, SQLite, RESTful API, 
MVC Mimarisi, Rol Tabanli Erisim Kontrolu

% ============================================================================
% 1.2 GIRIS
% ============================================================================
\section{Giris}
\label{sec:giris}

Yirmi birinci yuzyilda bilgi teknolojilerinin hizla gelismesi, kurumsal 
sureclerin dijitallestirilmesini zorunlu hale getirmistir. Bu dijital 
donusum, organizasyonlarin operasyonel verimliligini artirirken, ayni 
zamanda kaynak yonetimini optimize etmekte ve karar alma sureclerini 
hizlandirmaktadir.

Universite kampusleri, genis fiziksel alanlari ve yogun insan trafigi ile 
temizlik hizmetlerinin yonetimi acisindan onemli zorluklar barindirmaktadir. 
Kirikkale Universitesi, bunyesindeki cok sayida fakulte, bina ve hizmet 
alani ile bu zorlugun somut bir ornegini temsil etmektedir. Geleneksel 
temizlik yonetimi yontemleri, genellikle kagit tabanli kayitlar, sozel 
iletisim ve manuel takip sureclerine dayanmakta olup, bu durum cesitli 
sorunlara yol acmaktadir.

\subsection{Mevcut Sistemin Analizi}
\label{subsec:mevcut-sistem}

Mevcut temizlik yonetim sureclerinde karsilasilan temel problemler su 
sekilde ozetlenebilir:

\begin{enumerate}
    \item \textbf{Izlenebilirlik Eksikligi:} Temizlik gorevlerinin ne zaman 
    ve kim tarafindan yapildiginin takibi oldukca guctur. Kagit tabanli 
    kayitlar kaybolabilmekte veya manipule edilebilmektedir.
    
    \item \textbf{Koordinasyon Zorluklari:} Farkli binalar ve fakulteler 
    arasindaki temizlik personelinin koordinasyonu, ozellikle acil 
    durumlarda ciddi sorunlara neden olmaktadir.
    
    \item \textbf{Kaynak Optimizasyonu:} Personel is yukuun dengeli 
    dagitilamamamasi, bazi alanlarda asiri yuklenme, bazi alanlarda ise 
    yetersiz hizmet sunumu ile sonuclanmaktadir.
    
    \item \textbf{Raporlama Guclukleri:} Temizlik performansina iliskin 
    istatistiksel verilerin toplanmasi ve analiz edilmesi manuel 
    sureclerle oldukca zaman alicidir.
    
    \item \textbf{Iletisim Kopukluklari:} Yonetim, denetleyiciler ve 
    temizlik personeli arasindaki iletisim cogunlukla anlik ve 
    yapilandirilmamis bir sekilde gerceklesmektedir.
\end{enumerate}

\subsection{Cozum Yaklasimi}
\label{subsec:cozum-yaklasimi}

Bu calismada, yukarida belirtilen problemlere cozum uretmek amaciyla 
web tabanli bir temizlik takip ve yonetim sistemi gelistirilmistir. 
Gelistirilen sistem, modern yazilim muhendisligi prensiplerini takip 
ederek asagidaki temel ozellikleri sunmaktadir:

\begin{itemize}
    \item \textbf{Merkezi Veri Yonetimi:} Tum temizlik verileri merkezi 
    bir veritabaninda saklanarak tutarlilik ve erisilebilirlik 
    saglanmaktadir.
    
    \item \textbf{Rol Tabanli Erisim Kontrolu (RBAC):} Dort farkli 
    kullanici rolu ile hiyerarsik bir yetkilendirme sistemi 
    uygulanmaktadir.
    
    \item \textbf{Gercek Zamanli Takip:} Gorevlerin anlik durumlari 
    sistem uzerinden izlenebilmektedir.
    
    \item \textbf{Otomatik Gorev Atama:} Haftalik programlara dayali 
    olarak gunluk gorevler otomatik olarak olusturulmaktadir.
    
    \item \textbf{Kapsamli Raporlama:} Fakulte, personel ve donem 
    bazinda detayli performans raporlari uretilebilmektedir.
    
    \item \textbf{Platformdan Bagimsiz Erisim:} PWA mimarisi sayesinde 
    sistem, masaustu ve mobil cihazlardan erisilebilir durumdadir.
\end{itemize}

\subsection{Projenin Kapsami}
\label{subsec:proje-kapsami}

Bu bitirme projesi, Kirikkale Universitesi Muhendislik Fakultesi'nden 
baslayarak pilot uygulama kapsaminda alti fakulteyi iceren bir 
temizlik takip sistemi gelistirmeyi hedeflemektedir. Sistemin 
kapsami asagidaki bilesenleri icermektedir:

\begin{enumerate}
    \item \textbf{Kullanici Yonetimi Modulu:} Kullanici kayit, giris, 
    oturum yonetimi ve rol atama islevleri.
    
    \item \textbf{Organizasyon Yapisi Modulu:} Fakulte, bina, departman 
    ve lokasyon (oda) hiyerarsisinin tanimlanmasi.
    
    \item \textbf{Program Yonetimi Modulu:} Haftalik temizlik 
    programlarinin olusturulmasi ve personel atamalarinin yapilmasi.
    
    \item \textbf{Gorev Yonetimi Modulu:} Gunluk gorevlerin otomatik 
    olusturulmasi, takibi ve onay sureclerinin yonetimi.
    
    \item \textbf{Tatil Yonetimi Modulu:} Resmi tatil ve ozel gunlerin 
    tanimlanmasi ve gorev sistemine entegrasyonu.
    
    \item \textbf{Raporlama Modulu:} Cesitli kriterlere gore performans 
    ve istatistik raporlarinin olusturulmasi.
\end{enumerate}

\subsection{Raporun Organizasyonu}
\label{subsec:rapor-organizasyonu}

Bu bitirme projesi raporu, akademik standartlara uygun olarak 
yapilandirilmis olup, asagidaki bolumlerden olusmaktadir:

\textbf{Bolum 2 - Problem Tanimi:} Bu bolumde, projenin ele aldigi 
problem detayli olarak tanimlanmakta ve mevcut durumun analizi 
yapilmaktadir. Ayrica projenin amac ve hedefleri net bir sekilde 
ortaya konulmaktadir.

\textbf{Bolum 3 - Benzer Calismalar:} Literaturde yer alan ilgili 
akademik calismalar ve mevcut ticari cozumler incelenmektedir. 
Bu bolumde, projenin mevcut calismalardan farkliliklari ve 
katkilari tartisilmaktadir.

\textbf{Bolum 4 - Materyal ve Yontem:} Projede kullanilan 
programlama dilleri, framework'ler, veritabani sistemleri ve 
gelistirme araclari detayli olarak aciklanmaktadir. Ayrica 
yazilim gelistirme metodolojisi ve mimari kararlar bu bolumde 
ele alinmaktadir.

\textbf{Bolum 5 - Gelistirilen Sistem:} Sistemin teknik detaylari, 
veritabani semasi, API tasarimi, kullanici arayuzu ve is akislari 
diyagramlar esliginde sunulmaktadir.

\textbf{Bolum 6 - Sonuc:} Projenin ciktilari degerlendirilmekte, 
karsilasilan zorluklar tartisilmakta ve gelecek calismalar icin 
oneriler sunulmaktadir.

\subsection{Projenin Onemi ve Katkilari}
\label{subsec:proje-onemi}

Bu proje, asagidaki acilardan onem tasimakta ve literature katki 
saglamaktadir:

\begin{enumerate}
    \item \textbf{Akademik Katki:} Kurumsal temizlik yonetimi 
    alaninda Turkce literature kapsamli bir uygulama calismasi 
    eklenmektedir.
    
    \item \textbf{Pratik Katki:} Kirikkale Universitesi'nin temizlik 
    operasyonlarinin dijitallestirilmesi ile somut bir verimlilik 
    artisi hedeflenmektedir.
    
    \item \textbf{Teknolojik Katki:} Modern web teknolojilerinin 
    kurumsal yonetim sistemlerinde etkin kullanimina ornek 
    teskil etmektedir.
    
    \item \textbf{Metodolojik Katki:} Rol tabanli erisim kontrolu 
    ve hiyerarsik organizasyon yonetimi konularinda referans 
    bir uygulama sunulmaktadir.
\end{enumerate}

% ============================================================================
% 1.3 TEMEL KAVRAMLAR
% ============================================================================
\section{Temel Kavramlar ve Terminoloji}
\label{sec:temel-kavramlar}

Bu bolumde, raporda kullanilan teknik terimlerin ve kavramlarin 
aciklamalari verilmektedir.

\subsection{Web Uygulama Kavramlari}
\label{subsec:web-kavramlari}

\begin{description}
    \item[Progressive Web Application (PWA):] Web ve mobil uygulama 
    ozelliklerini birlestiren, tarayici uzerinden calisan ancak 
    yerel uygulama deneyimi sunan modern web uygulamasi mimarisidir. 
    PWA'lar, cevrimdisi calisabilme, bildirim gonderme ve ana ekrana 
    eklenebilme gibi ozellikler sunar.
    
    \item[Single Page Application (SPA):] Kullanici etkileşimlerinde 
    sayfanin tamaminin yeniden yuklenmesi yerine, yalnizca gerekli 
    bilesenlerin dinamik olarak guncellenmesi prensibine dayanan 
    web uygulama mimarisidir.
    
    \item[RESTful API:] Representational State Transfer (REST) 
    prensiplerine uygun olarak tasarlanmis, HTTP protokolu uzerinden 
    kaynak odakli iletisim saglayan uygulama programlama arayuzudur.
    
    \item[Model-View-Controller (MVC):] Yazilim uygulamalarini uc 
    temel bilesene ayiran mimari desendir: Model (veri ve is 
    mantigi), View (kullanici arayuzu) ve Controller (kullanici 
    girislerinin islenmesi).
\end{description}

\subsection{Guvenlik Kavramlari}
\label{subsec:guvenlik-kavramlari}

\begin{description}
    \item[Rol Tabanli Erisim Kontrolu (RBAC):] Kullanicilara belirli 
    roller atanarak, bu rollere gore sistem kaynaklarina erisim 
    yetkilerinin belirlenmesi yaklasimdir.
    
    \item[Oturum Yonetimi (Session Management):] Kullanici 
    kimlik dogrulamasi sonrasi olusturulan oturum bilgilerinin 
    sunucu tarafinda saklanmasi ve yonetilmesi surecidir.
    
    \item[Sifre Hashleme:] Kullanici sifrelerinin geri donusturulemez 
    kriptografik fonksiyonlar araciligiyla guvenli bir sekilde 
    saklanmasi islemidir. Bu projede bcrypt algoritmasi 
    kullanilmaktadir.
    
    \item[CORS (Cross-Origin Resource Sharing):] Farkli kaynaklardan 
    gelen HTTP isteklerinin guvenli bir sekilde yonetilmesini 
    saglayan mekanizmadir.
\end{description}

\subsection{Veritabani Kavramlari}
\label{subsec:veritabani-kavramlari}

\begin{description}
    \item[Iliskisel Veritabani:] Verilerin tablolar halinde organize 
    edildigi ve tablolar arasinda iliskilerin tanimlandigi 
    veritabani modelidir.
    
    \item[SQLite:] Sunucu gerektirmeyen, dosya tabanli, hafif ve 
    tasinabilir bir iliskisel veritabani yonetim sistemidir.
    
    \item[Foreign Key (Yabanci Anahtar):] Bir tablodaki alana, 
    baska bir tablonun birincil anahtarina referans vererek 
    tablolar arasi iliski kuran kisitlamadir.
    
    \item[CRUD Islemleri:] Create (Olustur), Read (Oku), Update 
    (Guncelle), Delete (Sil) temel veritabani operasyonlarinin 
    kisaltmasidir.
\end{description}

% ============================================================================
% 1.4 SISTEM GEREKSINIMLERI
% ============================================================================
\section{Sistem Gereksinimleri}
\label{sec:sistem-gereksinimleri}

Bu bolumde, KTTYS'nin fonksiyonel ve fonksiyonel olmayan 
gereksinimleri tanimlanmaktadir.

\subsection{Fonksiyonel Gereksinimler}
\label{subsec:fonksiyonel-gereksinimler}

Sistemin karsilamasi gereken temel fonksiyonel gereksinimler 
Tablo \ref{tab:fonksiyonel-gereksinimler}'de listelenmistir.

\begin{longtable}{|p{1.5cm}|p{3cm}|p{8cm}|}
\hline
\textbf{Kod} & \textbf{Gereksinim} & \textbf{Aciklama} \\
\hline
\endfirsthead
\hline
\textbf{Kod} & \textbf{Gereksinim} & \textbf{Aciklama} \\
\hline
\endhead
\hline
\endfoot

FR-01 & Kullanici Girisi & Sistem, kullanicilarin kullanici adi ve 
sifre ile guvenli bir sekilde giris yapabilmesini saglamalidir. \\
\hline
FR-02 & Rol Yonetimi & Sistem, dort farkli kullanici rolunu (Super 
Admin, Fakulte Admin, Supervisor, Staff) desteklemelidir. \\
\hline
FR-03 & Fakulte Yonetimi & Super Admin, fakulte ekleme, duzenleme 
ve silme islemlerini yapabilmelidir. \\
\hline
FR-04 & Kullanici Yonetimi & Yetkili kullanicilar, personel ekleme, 
duzenleme ve pasiflestirme islemlerini yapabilmelidir. \\
\hline
FR-05 & Lokasyon Yonetimi & Bina, departman ve oda tanimlamalari 
yapilabilmelidir. \\
\hline
FR-06 & Program Olusturma & Haftalik temizlik programlari 
olusturulabilmeli ve personele atanabilmelidir. \\
\hline
FR-07 & Gorev Yonetimi & Gunluk gorevler otomatik olusturulmali 
ve personel tarafindan takip edilebilmelidir. \\
\hline
FR-08 & Gorev Onaylama & Supervisor, tamamlanan gorevleri onaylama 
veya reddetme yetkisine sahip olmalidir. \\
\hline
FR-09 & Tatil Yonetimi & Resmi tatiller tanimlanabilmeli ve 
tatil gunlerinde gorev atamasi engellenmelidir. \\
\hline
FR-10 & Raporlama & Fakulte ve personel bazinda performans 
raporlari olusturulabilmelidir. \\
\hline

\caption{Sistemin Fonksiyonel Gereksinimleri}
\label{tab:fonksiyonel-gereksinimler}
\end{longtable}

\subsection{Fonksiyonel Olmayan Gereksinimler}
\label{subsec:fonksiyonel-olmayan-gereksinimler}

Sistemin kalite ozelliklerine iliskin gereksinimler asagida 
detaylandirilmistir:

\begin{enumerate}
    \item \textbf{Performans:}
    \begin{itemize}
        \item Sayfa yuklenme suresi 3 saniyeyi gecmemelidir.
        \item API yanit sureleri ortalama 500ms altinda olmalidir.
        \item Sistem, es zamanli 100 kullaniciyi desteklemelidir.
    \end{itemize}
    
    \item \textbf{Guvenlik:}
    \begin{itemize}
        \item Tum sifreler bcrypt ile hashlenmelidir.
        \item Oturum sureleri 24 saat ile sinirlandirilmalidir.
        \item Yetkisiz erisim girisimleri loglanmalidir.
    \end{itemize}
    
    \item \textbf{Kullanilabilirlik:}
    \begin{itemize}
        \item Arayuz Turkce olmalidir.
        \item Responsive tasarim ile mobil uyumluluk saglanmalidir.
        \item Kullanici hatalari anlasilir mesajlarla bildirilmelidir.
    \end{itemize}
    
    \item \textbf{Bakim Kolayligi:}
    \begin{itemize}
        \item Kod, moduler yapida organize edilmelidir.
        \item API endpoint'leri RESTful standartlarina uymalidir.
        \item Veritabani semasi genisletilebilir olmalidir.
    \end{itemize}
    
    \item \textbf{Tasinabilirlik:}
    \begin{itemize}
        \item Sistem, farkli isletim sistemlerinde calisabilmelidir.
        \item Modern web tarayicilarinin tamaminda uyumlu olmalidir.
    \end{itemize}
\end{enumerate}

% ============================================================================
% 1.5 PROJENIN SINIRLILIKLARI
% ============================================================================
\section{Projenin Sinirliliklari}
\label{sec:sinirliliklar}

Bu calismanin kapsami ve uygulama alanina iliskin bazi 
sinirliliklar bulunmaktadir:

\begin{enumerate}
    \item \textbf{Pilot Uygulama Kapsami:} Sistem, Kirikkale 
    Universitesi'nin mevcut alti fakultesi icin tasarlanmis 
    olup, hastane ve diger kampus binalari kapsam disindadir.
    
    \item \textbf{Cevrimdisi Calisma:} PWA mimarisine ragmen, 
    tam cevrimdisi calisma destegi bu asamada sinirlidir.
    
    \item \textbf{Entegrasyon:} Mevcut universite bilgi sistemleri 
    ile entegrasyon bu proje kapsaminda gerceklestirilmemistir.
    
    \item \textbf{Mobil Uygulama:} Yerel mobil uygulama (native app) 
    gelistirmesi yapilmamis olup, mobil erisim PWA uzerinden 
    saglanmaktadir.
    
    \item \textbf{Coklu Dil Destegi:} Sistem yalnizca Turkce 
    arayuz sunmaktadir.
\end{enumerate}

Bu sinirliliklar, projenin gelecek asamalarinda ele alinmasi 
planlanan gelistirme alanlarini olusturmaktadir.
